% =====================================================================
%  Цель и задание лабораторной работы.
% =====================================================================

\textbf{Цель:} освоить на практике основные принципы создания систем
анализа и распознавания речи.

\vspace{1em}

\textbf{Задание (вариант 1)}: разработать автоматизированную систему
анализа (распознавания) речи со следующей комбинацией признаков
варианта:
\begin{itemize}
  \item \textbf{Поддерживаемые языки} -- английский;
  \item \textbf{Предметная область} -- научные статьи по
  computer science.
\end{itemize}

\vspace{1em}

Система анализа должна обеспечивать следующие минимальные
возможности:
\begin{itemize}
  \item задание списка операций, на которые система может реагировать;
  \item организацию автоматической реакции системы на речевой сигнал с
  уведомлением пользователя о происходящем;
  \item настройки (выбор) естественного языка.
\end{itemize}

В отчёт включены: описание структуры разработанной системы,
технологий и методик; основные алгоритмы и принципы реализации
компонент системы (блок-схемы алгоритмов); результаты тестирования;
результаты анализа полученных данных и предложения по улучшению
работы системы.

\vspace{1em}

Распознавание речи реализовано как продолжение ЛР8 и как ещё один
способ управлять уже существующей информационно-поисковой системой
(ЛР1--ЛР4): голосом можно продиктовать поисковый запрос, перейти в
любой раздел, озвучить открытый документ научной статьи или остановить
озвучивание. Распознавание выполняет локальная нейросетевая модель
Whisper (библиотека faster-whisper), без обращения к облачным
сервисам (см. раздел «Технологии и методики»). Основной язык
распознавания -- английский (вариант~1); дополнительно, сверх
обязательного по варианту, поддержан русский язык.
